\chapter{Liver Biopsy}

\section*{Overview}
A liver biopsy is a procedure that obtains a small core of liver tissue, usually percutaneously, to diagnose and stage liver disease.

\section*{Alternative Names}
Percutaneous liver biopsy; liver needle biopsy; hepatic biopsy

\section*{Principle}
A spring-loaded or suction biopsy needle is advanced into the liver through the intercostal space or via a transjugular route, capturing a tissue core for histologic evaluation of inflammation, fibrosis, steatosis, and other pathology.

\section*{History}
Percutaneous liver biopsy was introduced by Iversen and Roholm in 1939 and became a standard diagnostic tool, increasingly guided by ultrasound in recent decades.

\section*{Why This Test or Procedure Is Ordered}
Performed for unexplained liver enzyme elevation, chronic hepatitis of uncertain cause, staging of fibrosis in chronic liver disease, evaluation of autoimmune hepatitis, and characterization of hepatic masses.

\section*{Is This a Primary or Secondary Test or Procedure}
Primary diagnostic procedure for parenchymal liver disease; secondary when non-invasive tests are inconclusive.

\section*{How This Test or Procedure Is Performed}
Performed with the patient supine, usually under ultrasound guidance. The biopsy site is marked and anesthetized, the needle is inserted during a breath-hold, and one or two cores are obtained. Pressure is applied, and the patient is observed.

\section*{What Preparation Is Needed}
Coagulation studies, platelet count, and blood type. Anticoagulants and antiplatelet agents are held, and informed consent is obtained. NPO is generally not required.

\section*{Contraindications and Precautions}
Coagulopathy, thrombocytopenia, ascites, or an uncooperative patient are relative contraindications. Precaution: a transjugular approach may be used when bleeding risk is high.

\section*{Risks}
Pain, bleeding with hemoperitoneum or subcapsular hematoma, biliary leak, and rarely puncture of other organs or death.

\section*{What Happens After the Test or Procedure}
The patient is observed for 2-4 hours with vital sign monitoring before discharge; severe pain or hypotension prompts urgent evaluation.

\section*{Understanding Your Results}
Histology grades necroinflammation, stages fibrosis, identifies the cause (steatohepatitis, autoimmune hepatitis, viral hepatitis, drug injury), and assesses response to therapy.

\section*{Normal Results}
Normal hepatocytes with no significant inflammation, fibrosis, steatosis, or other abnormality.

\section*{Follow-up Tests and Procedures}
Results guide treatment; serial non-invasive fibrosis markers or repeat biopsy monitor disease progression.

\section*{References}
\begin{enumerate}
  \item MedlinePlus. Liver biopsy. U.S. National Library of Medicine; 2025.
  \item Current Medical Diagnosis and Treatment. McGraw-Hill; 2025.
\end{enumerate}

\clearpage
